% Generated draft: requires mathematical and source review.
\documentclass[11pt]{article}
\usepackage[margin=25mm]{geometry}
\usepackage{fontspec}
\setmainfont{Latin Modern Roman}
\usepackage{amsmath,amssymb,mathtools,unicode-math}
\setmathfont{Latin Modern Math}
\usepackage{enumitem,needspace,xurl,hyperref}
\setlength{\parindent}{0pt}
\setlength{\parskip}{5pt}
\setlength{\emergencystretch}{3em}
\newcommand{\N}{\mathbb{N}}
\newcommand{\Z}{\mathbb{Z}}
\newcommand{\Q}{\mathbb{Q}}
\newcommand{\R}{\mathbb{R}}
\newcommand{\C}{\mathbb{C}}
\newcommand{\F}{\mathbb{F}}
\newcommand{\A}{\mathbb{A}}
\newcommand{\PP}{\mathbb{P}}
\newcommand{\E}{\mathbb{E}}
\newcommand{\eps}{\varepsilon}
\newcommand{\dd}{\,\mathrm{d}}
\begin{document}
\section*{UnsolvedMath ID: 800019}
\textbf{Source problem code:} catalogue entry
\subsection{Definitions and mathematical statement}
A semigroup is a nonempty set with an associative multiplication. It is commutative if \(xy=yx\) for all elements \(x,y\). It is finitely generated if there is a finite subset \(A\) such that every element is a product of a positive finite number of elements of \(A\). A commutative semigroup \(S\) is separative if
\[
\forall x,y\in S,\qquad x^2=xy=y^2\ \Longrightarrow\ x=y.
\]
\textbf{Problem (A. A. Borisov).} Describe the finitely generated commutative separative \(E\)-semigroups.

\textbf{Unresolved qualification.} The supplied record does not define an \(E\)-semigroup. Thus the additional property denoted by \(E\) must be recovered from the source before this classification request has a fully specified mathematical domain. No definition of that property is assumed here, and the accompanying endomorphism statement alone does not determine it.

\textbf{Supplied background assertion.} For a commutative semigroup \(S\), let \(\operatorname{End}(S)\) be its set of semigroup endomorphisms, equipped with pointwise multiplication
\[
(fg)(s)=f(s)g(s)\qquad(f,g\in\operatorname{End}(S),\ s\in S).
\]
This multiplication makes \(\operatorname{End}(S)\) a commutative semigroup. Under the usual convention, a commutative semigroup \(T\) is Archimedean if
\[
\forall a,b\in T\ \exists n\in\Z_{\geq 1}\ \exists u\in T^1:\quad a^n=bu,
\]
where \(T^1\) denotes \(T\) with an identity adjoined when necessary. The source reports, with reference label \([\mathrm{Borisov4}]\), that for a finitely defined commutative semigroup it is algorithmically decidable whether its endomorphism semigroup under pointwise multiplication is Archimedean. Here ``finitely defined'' is read in its usual sense of admitting a finite presentation by generators and relations. This reported decidability assertion is background, not a resolution of the classification request.
\subsection{Short English statement}
The problem, attributed to A. A. Borisov, asks for a description of finitely generated commutative separative \(E\)-semigroups. The meaning of \(E\)-semigroup is missing from the supplied record, so the target class remains incompletely specified. The source also reports a decidability result about Archimedean endomorphism semigroups for finitely presented commutative semigroups.
\subsection{Sources}
\begin{itemize}
\item Local UnsolvedMath registry, permanent ID \(800019\), source code \texttt{catalogue entry}, title ``The Lyapin's Notebook --- Problem 19.'' The supplied registry labels the problem open; this status has not been independently verified.
\item Bershadsky, Kublanovsky, and Mashevitzky, \emph{The Lyapin's Notebook: Unsolved Problems in Semigroup Theory} (2026), Problem 19. Supplied source URL: \url{https://arxiv.org/abs/2604.04763}. The problem text attributes the question to A. A. Borisov; the registry's proposer fields name the three compilation authors and give the year 2026.
\item Supplied provenance records access on 2026-10-16 and extraction by \texttt{arxiv\_latex}. These are supplied provenance details, not a claim of independent source consultation or verification.
\item Reference label \([\mathrm{Borisov4}]\) is supplied for the background decidability assertion. Its bibliographic entry is absent from the record, so that citation cannot be identified precisely here.
\end{itemize}
\end{document}
